\begin{vidu} %[2P4-6.2K][Loại 2]
	Hai dây dẫn thẳng dài vô hạn song song đặt cách nhau 10 cm, dòng điện chạy qua hai dây là $I_1 = I_2 = 10$ A cùng chiều. Lực tác dụng lên mỗi mét chiều dài của dây là
	\choice
	{lực hút có độ lớn $2. 10^{-7}$ N}
	{lực đẩy có độ lớn $2. 10^{-7}$ N}
	{\True lực hút có độ lớn $2. 10^{-4}$ N}
	{lực đẩy có độ lớn $2. 10^{-4}$ N}
	\loigiai{\begin{itemize}
			\item Cảm ứng từ do dây 1 tương tác lên dây 2 là: $B=2. 10^{-7}\dfrac{I_1}{r}$
			\item Lực tương tác 1 m dài dây 2 lực:
			$F=2. 10^{-7}\dfrac{I_1I_2}{r}$
			\item Tương tự, dây 2 cũng tương tác lên 1 m dài dây 1 lực có độ lớn như vậy.
			\item Áp dụng quy tắc bàn tay trái $\Rightarrow$ hai dây có dòng điện cùng chiều thì hút nhau, ngược chiều thì đẩy nhau.
			\item Hai dây dẫn mang dòng điện cùng chiều $\Rightarrow$ hai dây hút nhau
			$F=2. 10^{-7}\dfrac{I_1I_2}{r}=2. 10^{-7}\dfrac{10^2}{0,1}=2.10^{-4}$ N
	\end{itemize}}
\end{vidu}
\begin{vidu} %[2P4-6.2K][Loại 2]
	Dây dẫn thẳng dài $I_1 = 10$ A, đặt dây dẫn khác song song cách $I_1$ một khoảng 10 cm có dòng $I_2 = 5$ A chạy ngược chiều dòng $I_1$. Lực từ tác dụng lên mỗi mét chiều dài dây mang dòng điện $I_2$ là
	\choice
	{lực hút $2. 10^{-4}$ N}
	{\True lực đẩy $10^{-4}$ N}
	{lực hút $10^{-4}$ N}
	{lực đẩy $2.10^{-4}$ N}
	\loigiai{Hai dây dẫn mang dòng điện ngược chiều $\Rightarrow$ hai dây đẩy nhau: $F=2. 10^{-7}\dfrac{I_1I_2}{r}=2. 10^{-7}\dfrac{10. 5}{0,1}=10^{-4}$ N.
	}
\end{vidu}
\begin{vidu} %[2P4-6.2G][Loại 1]
	immini[thm]{Ba dòng điện thẳng song song $I_1, I_2$ và $I_3$ ở trong cùng một mặt phẳng, cho $I_1 = 20$ A, $I_2 = 15$ A, $I_3 = 25$ A. Khoảng cách giữa $I_1$ và $I_2$ là $a = 5$ cm, giữa $I_2$ và $I_3$ là $b = 3$ cm. Lực tác dụng lên 1 m chiều dài của $I_2$ là
		\choice
	{$3,7.10^{-5}$ N}
	{$12.10^{-4}$ N}
	{\True $37.10^{-4}$ N}
	{$25.10^{-4}$ N}}
{\begin{tikzpicture}[scale=1,thick,>=stealth']
		\coordinate (a1) at (0,0);
		\draw [very thick,dashed](a1)--++(0:0.5)coordinate (b1);
		\draw [very thick,-<](b1)--++(0:1.5)coordinate (c1)node[below=3pt]{\normalsize{\bth{I_{1}}}};
		\draw [very thick](b1)--++(0:3)coordinate (d1);
		\draw [very thick,dashed](d1)--++(0:0.5);
		\coordinate (a2) at (0,-2);
		\draw [very thick,dashed](a2)--++(0:0.5)coordinate (b2);
		\draw [very thick,->](b2)--++(0:1.5)coordinate (c2)node[above=3pt]{\normalsize{\bth{I_{2}}}};
		\draw [very thick](b2)--++(0:3)coordinate (d2);
		\draw [very thick,dashed](d2)--++(0:0.5);
		\coordinate (a3) at (0,-3.2);
		\draw [very thick,dashed](a3)--++(0:0.5)coordinate (b3);
		\draw [very thick,->](b3)--++(0:1.5)coordinate (c3)node[below=3pt]{\normalsize{\bth{I_{3}}}};
		\draw [very thick](b3)--++(0:3)coordinate (d3);
		\draw [very thick,dashed](d3)--++(0:0.5);
		\draw[<->] (1,0)--node[midway,fill=white]{\normalsize{5 cm}}(1,-2);
		\draw[<->] (1,-2)--node[midway,fill=white]{\normalsize{3 cm}}(1,-3.2);
		\end{tikzpicture}}
	\loigiai{
	immini[thm]{\begin{itemize}
			\item Lực từ do dòng $I_1$ tác dụng lên 1 m dòng $I_2$ là $F_{12} = 2.10^{-7}\dfrac{I_1I_2}{r_1} = 2.10^{-7}\dfrac{20.15}{0,05} = 1,2.10^{-3}$ N.
			\item Lực từ dò dòng $I_3$ tác dụng lên 1 m dòng $I_2$ là $F_{32} = 2.10^{-7}\dfrac{I_2I_3}{r_2} = 2.10^{-7}\dfrac{15.25}{0,03} = 2,5.10^{-3}$ N.
			\item Lực tổng hợp F tác dụng lên mỗi mét dây dẫn của dòng điện $I_2$ là $F=F_{12}+F_{32}=3,7.10^{-3}$ N.
	\end{itemize}}{\begin{tikzpicture}[scale=1,thick,>=stealth']
\coordinate (a1) at (0,0);
\draw [very thick,dashed](a1)--++(0:0.5)coordinate (b1);
\draw [very thick,-<](b1)--++(0:1.5)coordinate (c1)node[below=3pt]{\normalsize{\bth{I_{1}}}};
\draw [very thick](b1)--++(0:3)coordinate (d1);
\draw [very thick,dashed](d1)--++(0:0.5);
\coordinate (a2) at (0,-2);
\draw [very thick,dashed](a2)--++(0:0.5)coordinate (b2);
\draw [very thick,->](b2)--++(0:1.5)coordinate (c2)node[above=3pt]{\normalsize{\bth{I_{2}}}};
\draw [very thick](b2)--++(0:3)coordinate (d2);
\draw [very thick,dashed](d2)--++(0:0.5);
\coordinate (a3) at (0,-3.2);
\draw [very thick,dashed](a3)--++(0:0.5)coordinate (b3);
\draw [very thick,->](b3)--++(0:1.5)coordinate (c3)node[below=3pt]{\normalsize{\bth{I_{3}}}};
\draw [very thick](b3)--++(0:3)coordinate (d3);
\draw [very thick,dashed](d3)--++(0:0.5);
\draw[<->] (1,0)--node[midway,fill=white]{\normalsize{5 cm}}(1,-2);
\draw[<->] (,-2)--node[midway,fill=white]{\normalsize{3 cm}}(1,-3.2);
\draw[very thick,->](2.5,-2)--++(-90:0.7)coordinate(f1)node[left]{$\vec{F}_{12}$};
\draw[very thick,->](2.6,-2)--++(-90:0.8)coordinate(f2)node[right]{$\vec{F}_{32}$};
\end{tikzpicture}
}}
\end{vidu}
\begin{vidu} %[2P4-6.2G][Loại 1] 
immini[thm]{Dòng điện thẳng dài đặt song song với nhau đi qua ba đỉnh của một tam giác theo phương vuông góc với mặt phẳng hình vẽ. Cho các dòng điện chạy qua có chiều như hình vẽ với các cường độ dòng điện $I_1 = 10$ A, $I_2 = 20$ A, $I_3 = 30$ A. Biết $I_1$ cách $I_2$ và $I_3$ lần lượt là 8 cm và 6 cm. Lực tổng hợp tác dụng lên mỗi mét dây dẫn có dòng điện $I_1$ là
\choice
{$2.10^{-3}$ N}
{$1,2.10^{-3}$ N}
{\True $1,12.10^{- 3}$ N}
{$1,5.10^{-3}$ N}}{\begin{tikzpicture}[scale=1,thick,>=stealth']
	\coordinate (a) at (0,0);
	\draw [dashed](a)--++(270:3)coordinate (b);
	\draw [dashed](b)--++(0:3)coordinate (c);
	\draw [dashed](a)--(b)--(c)--cycle;
	\draw[fill=white] (a) circle(0.2)node[left=5pt]{\normalsize{\bth{I_{2}}}};
	\filldraw[black] (a) circle(0.05);
	\draw[fill=white] (b) circle(0.2)node[left=5pt]{\normalsize{\bth{I_{1}}}};
	\filldraw[black] (b) circle(0.05);
	\draw[fill=white] (c) circle(0.2)node[right=5pt]{\normalsize{\bth{I_{3}}}};
	\filldraw[black] (c) circle(0.05);
	\end{tikzpicture}
}
\loigiai{
immini[thm]{\begin{itemize}
\item Lực từ do dòng $I_2$ tác dụng lên 1 m dòng $I_1$ là $F_{21} = 2.10^{-7}\dfrac{I_1I_2}{r_1} = 2.10^{-7}\dfrac{20.10}{0,08} = 5.10^{-4}$ N.
\item Lực từ dò dòng $I_3$ tác dụng lên 1 m dòng $I_1$ là $F_{31} = 2.10^{-7}\dfrac{I_3I_1}{r_2} = 2.10^{-7}\dfrac{30.10}{0,06} = 10^{-3}$ N.
\item Lực tổng hợp F tác dụng lên mỗi mét dây dẫn của dòng điện $I_1$ là $F = \sqrt {F_{21}^2 + F_{31}^2} = 1,12.10^{-3}$ N.
\end{itemize}}{\begin{tikzpicture}[scale=1,thick,>=stealth']
\coordinate (a) at (0,0);
\draw [dashed](a)--++(270:3)coordinate (b);
\draw [dashed](b)--++(0:3)coordinate (c);
\draw [dashed](a)--(b)--(c)--cycle;
\draw[very thick,->](b)--++(90:1)coordinate(f1)node[left]{$\vec{F}_{21}$};
\draw[very thick,->](b)--++(0:1)coordinate(f2)node[below]{$\vec{F}_{31}$};
\path(f2)--++(90:1)coordinate(f);
\draw[very thick,->](b)--(f)node[above]{$\vec{F}$};
\draw [dashed](f1)--(f)(f2)--(f);

\draw[fill=white] (a) circle(0.2)node[left=5pt]{\normalsize{\bth{I_{2}}}};
\filldraw[black] (a) circle(0.05);
\draw[fill=white] (b) circle(0.2)node[left=5pt]{\normalsize{\bth{I_{1}}}};
\filldraw[black] (b) circle(0.05);
\draw[fill=white] (c) circle(0.2)node[right=5pt]{\normalsize{\bth{I_{3}}}};
\filldraw[black] (c) circle(0.05);
\end{tikzpicture}}}
\end{vidu}
\begin{vidu} %[2P4-6.2G][Loại 1]
	immini[thm]{Ba dòng điện thẳng dài đặt song song với nhau, cách đều nhau đi qua ba đỉnh của một tam giác đều cạnh 4 cm theo phương vuông góc với mặt phẳng hình vẽ. Biết các cường độ dòng điện $I_1 = 10$ A, $I_2=I_3= 20$ A. Lực tổng hợp F tác dụng lên mỗi mét dây dẫn có dòng điện $I_1$ là
		\choice
	{$2,5.10^{-3}$ N}
	{\True $1,73.10^{-3}$ N}
	{$2.10^{-3}$ N}
	{$4.10^{-3}$ N}}
{\begin{tikzpicture}[scale=1,thick,>=stealth']
		\coordinate (a) at (0,0)node[above=5pt]{\normalsize{\bth{A}}};
		\filldraw[black] (a) circle(0.05);
		\draw (a) circle(0.2)node[left=5pt]{\normalsize{\bth{I_{1}}}};
		
		\draw [dashed](a)--++(240:3)coordinate (b)node[below left=5pt]{\normalsize{\bth{B}}};
		\draw [dashed](b)--++(0:3)coordinate (c)node[below right=5pt]{\normalsize{\bth{C}}};
		
		\draw [dashed](a)--(b)--(c)--cycle;
		
		\filldraw[black] (b) circle(0.05);
		\draw (b) circle(0.2)node[above left=5pt]{\normalsize{\bth{I_{2}}}};
		\filldraw[black] (c) circle(0.05);
		\draw (c) circle(0.2)node[above right=5pt]{\normalsize{\bth{I_{3}}}};
		\end{tikzpicture}
	}
	\loigiai{
	immini[thm]{\begin{itemize}
			\item Lực từ do dòng $I_2$ tác dụng lên 1 m dòng $I_1$ là $F_{21} = 2.10^{-7}\dfrac{I_1I_2}{a} = 2.10^{-7}\dfrac{20.10}{0,04} = 10^{-3}$ N.
			\item Lực từ do dòng $I_3$ tác dụng lên 1 m dòng $I_1$ là $F_{31} = 2.10^{-7}\dfrac{I_3I_1}{a} = 2.10^{-7}\dfrac{20.10}{0,04} = 10^{-3}$ N.
			\item Lực tổng hợp F tác dụng lên mỗi mét dây dẫn của dòng điện $I_1$ là\\ $F = \sqrt {F_{21}^2 + F_{31}^2 + 2.F_{21}.F_{31}.\cos60^\circ} = {1,73.10^{- 3}}$ N.
	\end{itemize}}{\begin{tikzpicture}[scale=1,thick,>=stealth']
\coordinate (a) at (0,0)node[above=5pt]{\normalsize{\bth{A}}};
\draw [dashed](a)--++(240:3)coordinate (b)node[below left=5pt]{\normalsize{\bth{B}}};
\draw [dashed](b)--++(0:3)coordinate (c)node[below right=5pt]{\normalsize{\bth{C}}};
\draw[very thick,->](a)--++(-120:1)coordinate(f1)node[left]{$\vec{F}_{21}$};
\draw[very thick,->](a)--++(-60:1)coordinate(f2)node[right]{$\vec{F}_{31}$};
\path(f2)--++(-120:1)coordinate(f);
\draw[very thick,->](a)--(f)node[below]{$\vec{F}$};
\draw [dashed](f1)--(f)(f2)--(f);
\draw[fill=white]  (a) circle(0.2)node[left=5pt]{\normalsize{\bth{I_{1}}}};
\filldraw[black] (a) circle(0.05);
\draw [dashed](a)--(b)--(c)--cycle;
\draw[fill=white] (b) circle(0.2)node[above left=5pt]{\normalsize{\bth{I_{2}}}};
\filldraw[black] (b) circle(0.05);
\draw[fill=white] (c) circle(0.2)node[above right=5pt]{\normalsize{\bth{I_{3}}}};
\filldraw[black] (c) circle(0.05);
\end{tikzpicture}
}}
\end{vidu}
\begin{vidu} %[2P4-6.2G][Loại 1]
	Ba dòng điện thẳng song song cùng chiều gồm $I_1 = I_2 = 500$ A và $I_3$ cùng nằm trong mặt phẳng nằm ngang vuông góc với mặt phẳng hình vẽ tại M, N và C. immini[thm]{Biết $\widehat{MCN} = 120^\circ$, C cách MN 2,5 cm. Dòng điện $I_3$ chạy trong dây dẫn bằng đồng có đường kính 1,5 mm, khối lượng riêng 8,9 $g/cm^3$. Lấy $g = 10\ m/s^2$. Dây $I_1$ và $I_2$ được giữ cố định. Dây mang dòng $I_3$ nằm cân bằng. Cường độ dòng điện $I_3$ bằng}{\begin{tikzpicture}[scale=1,thick,>=stealth']
		\coordinate (a) at (0,0)node[above=5pt]{\normalsize{\bth{C}}};
		\draw [dashed](a)--++(150:3)coordinate (b)node[below left=5pt]{\normalsize{\bth{M}}};
		\draw [dashed](a)--++(30:3)coordinate (c)node[below right=5pt]{\normalsize{\bth{N}}};
		\draw [dashed](b)--(c);
		\draw [fill=white](a) circle(0.2)node[left=5pt]{\normalsize{\bth{I_{3}}}};
		\filldraw[black] (a) circle(0.05);
		\draw [fill=white](b) circle(0.2)node[above left=5pt]{\normalsize{\bth{I_{1}}}};
		\filldraw[black] (b) circle(0.05);
		\draw [fill=white](c) circle(0.2)node[above right=5pt]{\normalsize{\bth{I_{2}}}};
		\filldraw[black] (c) circle(0.05);
		\end{tikzpicture}}
	\choice
	{58,6 A}
	{68,6 A}
	{\True 78,6 A}
	{88,6 A}
	\loigiai{$\vec{F}_{13}+\vec{F}_{23}+\vec{P}=\vec{0}\Rightarrow \underbrace{\vec{F}_{13}+\vec{F}_{23}}_{\vec{F}_{3}}=-\vec{P} \Rightarrow F_3 = P$.\\
		Dễ thấy $P = mg = DSg, MC = 5$ cm $\Rightarrow I_3 = 78,6$ A}
\end{vidu}